\documentclass[11pt,a4paper]{article}
\usepackage[margin=25mm]{geometry}
\usepackage{amsmath,amssymb,amsthm}
\newtheorem{theorem}{Theorem}
\newcommand{\gen}[1]{\langle#1\rangle}
\title{Universal regular amalgamation from a witness}
\author{Human Compression project}\date{}
\begin{document}\maketitle
We recover the conditional existence assertion in Section~4 of the fixed
Schulte--Weiss survey. Its premise is that an amalgamation already exists;
the argument does not decide which local pairs have an amalgamation.
We retain the standard end criterion for a string C-group and the
correspondence between such marked groups and regular polytopes.
Both are the same attributed inputs used by our acquired chamber theory.
The construction below is standard universal-presentation reasoning,
written out to keep its existence premise visible.

\begin{theorem}
Let $K,L$ be regular polytopes of rank $n\geq2$. Suppose a polytope $P$
of rank $n+1$ has all facets isomorphic to $K$ and all vertex figures
isomorphic to $L$. Then there is a regular polytope $U$ with those local
types which covers every polytope with those local types, after choosing
base flags. Its group is obtained from the universal string group by
imposing all the marked facet and vertex-figure relations.
\end{theorem}
\begin{proof}
Let
\[
 W=\gen{s_0,\ldots,s_n\mid s_i^2=1,
                 (s_is_j)^2=1\ (|i-j|>1)}.
\]
For $I\subseteq\{0,\ldots,n\}$ write $W_I=\gen{s_i:i\in I}$.
Put $A=\{0,\ldots,n-1\}$, $B=\{1,\ldots,n\}$ and $C=A\cap B$.
The standard markings give onto homomorphisms
$W_A\to\operatorname{Aut}(K)$ and
$W_B\to\operatorname{Aut}(L)$, with kernels $N_A,N_B$.
Let $N$ be their normal closure in $W$, and put $G=W/N$.
Use $G_I$ for the subgroup on the indicated images of the marks.

The generators of $W$ act by the corresponding adjacency operations on
the flags of $P$. These operations are involutions, and nonconsecutive
ones commute, so this is an action of $W$ (with words applied in their
written order). Every relation in $N_A$ acts trivially on all flags:
it stays inside a facet, where the adjacency action is the regular
flag action of a copy of $K$. This assertion holds in every facet,
since an isomorphism of regular polytopes can send the chosen flag to
any flag. The same argument in vertex figures applies to $N_B$.
Consequently the action factors through $G$.

Fix a base flag $\Phi$ of $P$. The action of $G_A$ on the flags containing
the base facet is precisely the regular flag action of $K$.
More explicitly, the quotient map induces
$W_A/N_A\twoheadrightarrow G_A$, while the restriction of the action
back to that facet has kernel exactly $N_A$ in $W_A$.
These two observations show that $G_A\cong W_A/N_A$ as marked groups.
In particular $G_A$ is a C-group and its action on this facet's flags
is free. Similarly $G_B\cong W_B/N_B$, and its action on the flags
containing the base vertex is free.

We claim $G_A\cap G_B=G_C$. If $g$ belongs to the intersection, then
$\Phi^g$ has both the base vertex and the base facet. The flags containing
these two faces form the flag set of their section. They are connected
by adjacencies of types $1,\ldots,n-1$. There is therefore $h\in G_C$
with $\Phi^h=\Phi^g$. Both $h$ and $g$ lie in $G_A$, whose action on
the facet is free, so $g=h$. The reverse inclusion is immediate.
The end criterion now proves that $G$ is a string C-group. Its marked
regular polytope $U$ has facets $K$ and vertex figures $L$ by the exact
end-group identifications just proved.

For any other polytope $Q$ with those local types, the same local
argument shows that $N$ acts trivially on its flags. Choose a base
flag $\Psi$. Identify the flags of $U$ with $G$ in such a way that
adjacency of type $i$ is right multiplication by $s_i$. The map
\[
                  g\longmapsto\Psi^g
\]
is well-defined, onto, and commutes with each adjacency operation.
It induces a rank-preserving incidence map on faces: flags containing
a fixed face of rank $i$ form an orbit under adjacencies of all types
except $i$, and their images therefore contain one fixed rank-$i$
face of $Q$. This gives a face map independent of the chosen incident
flag. Images of a flag remain a flag, and adjacent flags remain
adjacent. Thus it is a polytope covering. This proves universality.
\end{proof}

The proof needs no regularity assumption on the witness $P$. In particular
it includes a chiral witness with regular facets and vertex figures, as
in the survey. It uses regularity of the two local types essentially:
their marked relations must act trivially on every local flag, and the
local action used to identify the middle intersection must be free.
The oriented amalgamation assertion for chiral constituents has different
input data, including relative enantiomorphic choices; it is retained
separately with its source attribution.

The group presentation describes the universal object when the premise
is satisfied. It does not show that every pair with matching intermediate
sections has a witness, does not supply an order algorithm, and does not
imply finiteness. Thus the theorem is the established conditional statement,
not a solution of the survey's amalgamation or finiteness questions.

\begin{thebibliography}{9}\small
\bibitem{survey} E. Schulte and A. Ivi\'c Weiss,
\emph{Problems on polytopes, their groups, and realizations},
Period. Math. Hungar. 53 (2006), 231--255; fixed arXiv math/0608397v1,
Section 4 and its attribution to P. McMullen and E. Schulte,
\emph{Abstract Regular Polytopes}. The standard end criterion and
marked reconstruction are retained inputs, not new results of this note.
\end{thebibliography}
\end{document}
